(1) 取 \(z=x+iy\)，复速度定义为 \(dW/dz=u_x-i u_y\)，选择
\[
W(z)=U\left(z+\frac{R^2}{z}\right)-\frac{i\Gamma}{2\pi}\log\frac zR,\qquad |z|\ge R.
\]
令 \(z=re^{i\vartheta}\)，对应流函数
\(\psi=U(r-R^2/r)\sin\vartheta-\Gamma(2\pi)^{-1}\ln(r/R)\)，因此在 \(r=R\) 上流函数恒定，\(u_r(R,\vartheta)=0\)。沿圆柱表面的速度为
\[
u_\vartheta(R,\vartheta)=-2U\sin\vartheta+\frac{\Gamma}{2\pi R}.
\]
远场 \(u_x\to U,\ u_y\to0\)，而
\(\oint\mathbf u\cdot d\boldsymbol\ell=\Gamma\) 是圆柱外任一环路的有向环量；对数复势只允许在流体域中有多值势，其速度和流函数可为单值。

(2) 因表面径向速度严格为零，停滞点满足
\[
\sin\vartheta_s=\frac{\Gamma}{4\pi UR}.
\]
当 \(|\Gamma|<4\pi UR\) 有两个不同的表面停滞点；当等号成立，两点在 \(\vartheta=\pi/2\) 或 \(3\pi/2\) 合并；当超过临界值时表面切向速度不再变号，停滞点移至圆柱外其他位置，不表示原有理想势流不存在。由同一流线的Bernoulli积分，
\[
p(R,\vartheta)=p_\infty+\frac{\rho}{2}
\left[U^2-\left(-2U\sin\vartheta+\frac{\Gamma}{2\pi R}\right)^2\right].
\]
记 \(c=\Gamma/(2\pi R)\)。表面最大压力在允许的 \(\sin\vartheta\in[-1,1]\) 上使 \(|c-2U\sin\vartheta|\) 最小，故
\(p_{\max}=p_\infty+\rho[U^2-(\max\{|c|-2U,0\})^2]/2\)。
最大速度绝对值为 \(|c|+2U\)，故
\(p_{\min}=p_\infty+\rho[U^2-(|c|+2U)^2]/2\)。

(3) 流体对圆柱单位长度的压力力为
\(d\mathbf F=-pR(\cos\vartheta,\sin\vartheta)d\vartheta\)。
将第二问平方展开，只有与 \(\sin\vartheta\) 成正比的交叉项能对 \(F_y\) 给非零贡献：
\(p=p_\infty+\rho U^2/2-2\rho U^2\sin^2\vartheta+2\rho Uc\sin\vartheta-\rho c^2/2\)。
因此
\[
F_x=-R\int_0^{2\pi}p\cos\vartheta\,d\vartheta=0,\qquad
F_y=-R\int_0^{2\pi}p\sin\vartheta\,d\vartheta=-2\pi\rho URc=-\rho U\Gamma.
\]
在所约定的逆时针正环量下，升力朝 \(-y\) 方向；更换环量符号会反转该方向。径向压力力处处经过轴心，所以总压力力矩严格为零。无黏势流仍能产生升力，是外部流场的横向动量预算结果；同时 \(F_x=0\) 是理想流的D'Alembert悖论。

(4) 完整速度场为
\[
u_r=U(1-R^2/r^2)\cos\vartheta,\qquad
u_\vartheta=-U(1+R^2/r^2)\sin\vartheta+\Gamma/(2\pi r).
\]
对同心任意 \(r>R\)，
\(\int_0^{2\pi}u_r r\,d\vartheta=0\) 和
\(\int_0^{2\pi}u_\vartheta r\,d\vartheta=\Gamma\)，分别核验无源质量连续性与环量守恒。圆柱表面不穿透，所以流体在固定壁面上的压力功率
\(-\int_{\partial S}p\,u_r\,d\ell=0\)；升力的存在与壁面不做功不矛盾。由于升力来自全局横向动量流，不能以没有尾流阻力来推出其为零；\(\Gamma\to0\) 时压力上下对称，升阻力都趋零。
